% 分野別類題: 定数係数2階線形非同次微分方程式（定数変化法）
% 未定係数法が使えない右辺（tan x, 1/cos x, e^x/x など）に対する定数変化法（ロンスキアン利用）の基本形 3 問。
% コンパイル: TEXINPUTS="../../../00_common:$TEXINPUTS" latexmk -xelatex problems.tex
\documentclass[a4paper,11pt]{article}
\usepackage{preamble}

\begin{document}

\kamokumei{類題演習 --- 定数係数2階線形非同次微分方程式（定数変化法）}

\shijibun{次の常微分方程式を解き、導出過程を示せ。右辺が未定係数法では扱えない形（$\tan x$、$1/\cos x$、$e^{x}/x$ など）であるため、対応する同次方程式の基本解 $y_{1}, y_{2}$ を用いた定数変化法（ロンスキアンの利用）により特殊解を求めること。}

\shomon{問1.}{次の常微分方程式の一般解を求めよ。}
\[
y'' + y = \tan x \qquad \left(-\frac{\pi}{2} < x < \frac{\pi}{2}\right)
\]

\shomon{問2.}{次の常微分方程式の一般解を求めよ。}
\[
y'' + y = \frac{1}{\cos x} \qquad \left(-\frac{\pi}{2} < x < \frac{\pi}{2}\right)
\]

\shomon{問3.}{次の初期値問題を解け。}
\[
y'' - 2y' + y = \frac{e^{x}}{x}, \qquad y(1) = e, \quad y'(1) = e \qquad (x > 0)
\]

\end{document}
